% Generated from tex/chapters/ch04/05_ソフトウェア構築の基礎.tex for the SWEBOK structure tree
\subsubsection{検証しやすく構築する}

検証しやすく構築するとは、コードを書いた人、テスト担当者、利用者が欠陥を見つけやすいようにソフトウェアを作ることである。これは、後でテストを追加すればよいという意味ではない。構築中から、テスト可能性を意識した設計と実装が必要である。

具体的には、単体テストを書きやすい関数分割、依存関係を差し替えやすい構造、複雑な言語機能の制限、ログによる挙動記録、コーディング標準、コードレビュー、入力値の検査、異常系の明示が役立つ。

検証しやすいコードでは、失敗が早く発見される。失敗が早く見つかるほど、原因に近い場所で修正できるため、修正費用を下げやすい。
